\documentclass{sample}
\usepackage{booktabs}
\usepackage{caption}
\begin{document}
\section{Greek letters}
Uppercase and lowercase Greek letters are a common symbology. These letters must be typed in math environments.

\begin{center}
	\begin{tabular}{llc}
	\toprule
	Name & Typeset \\
	\midrule
	alpha & $\alpha$\\
	beta & $\beta$\\
	gamma & $\gamma$\\
	delta & $\delta$\\
	epsilon & $\epsilon$\\
	zeta & $\zeta$\\
	eta & $\eta$\\
	theta & $\theta$\\
	iota & $\iota$\\
	kappa & $\kappa$\\
	lambda & $\lambda$\\
	mu & $\mu$\\
	nu & $\nu$\\
	xi & $\xi$\\
	pi & $\pi$\\
	rho & $\rho$\\
	sigma & $\sigma$\\
	tau & $\tau$\\
	upsilon & $\upsilon$\\
	phi & $\phi$\\
	chi & $\chi$\\
	psi & $\psi$\\
	omega & $\omega$\\
	varepsilon & $\varepsilon$\\
	vartheta & $\vartheta$\\
	varpi & $\varpi$\\
	varrho & $\varrho$\\
	varsigma & $\varsigma$\\
	varphi & $\varphi$\\
	digamma & $\digamma$\\
	varkappa & $\varkappa$\\
	\bottomrule
	\end{tabular}
	\captionof{table}{Lowercase Greek letters (see source for type or command)}
	\label{Ta:lowgreek}
\end{center}

\begin{center}
	\begin{tabular}{llc}
	\toprule
	Name & Typeset & Typeset \\
	\midrule
	alpha & A (just type A)\\
	gamma & $\Gamma$\\
	delta & $\Delta$\\
	theta & $\Theta$\\
	lambda & $\Lambda$\\
	xi & $\Xi$\\
	pi & $\Pi$\\
	sigma & $\Sigma$\\
	upsilon & $\Upsilon$\\
	phi & $\Phi$\\
	psi & $\Psi$\\
	omega & $\Omega$\\
	vargamma & $\varGamma$\\
	vardelta & $\varDelta$\\
	vartheta & $\varTheta$\\
	varlambda & $\varLambda$\\
	varxi & $\varXi$\\
	varpi & $\varPi$\\
	varsigma & $\varSigma$\\
	varupsilon & $\varUpsilon$\\
	varphi & $\varPhi$\\
	varpsi & $\varPsi$\\
	varomega & $\varOmega$\\
	\bottomrule
	\end{tabular}
	\captionof{table}{Uppercase Greek letters (see source for type or command)}
	\label{Ta:uppgreek2}
\end{center}
\newpage
\section{Delimiters}
Delimiters enclose expressions.Like Greek letters, delimiters always go inside math environments.
\begin{center}
	\begin{tabular}{llc}
	\toprule
	Name & Typeset \\
	\midrule
	left parenthesis & $($\\
	right parenthesis & $)$\\
	left bracket & $[$\\
	right bracket & $]$\\
	left brace & $\{$\\
	right brace & $\}$\\
	backslash & $\backslash$\\
	forward slash & $/$\\
	left angle bracket & $\langle$\\
	right angle bracket & $\rangle$\\
	vertical line & $|$\\
	double vertical line (note this one) & $\|$\\
	left floor & $\lfloor$\\
	right floor & $\rfloor$\\
	left ceiling & $\lceil$\\
	right ceiling & $\rceil$\\
	upward & $\uparrow$\\
	double upward & $\Uparrow$\\
	downward & $\downarrow$\\
	double downward & $\Downarrow$\\
	up-and-down & $\updownarrow$\\
	double up-and-down & $\Updownarrow$\\
	upper-left corner & $\ulcorner$\\
	upper-right corner & $\urcorner$\\
	lower-left corner & $\llcorner$\\
	lower-right corner & $\lrcorner$\\
	\bottomrule
	\end{tabular}
	\captionof{table}{Delimiters (see source for type or command)}
	\label{Ta:delimiters}
\end{center}
With the exception of the corner delimiters, it is possible to stretch delimiters with \textbackslash left and \textbackslash right, in the form \textbackslash left \textit{delim1} \textbackslash right \textit{delim2}.

For example:
\[
\left|\frac{a+b}{2}\right|
\]

Note that the slightly modified command \textbackslash left. introduces a blank delimiter (as does the right counterpart) and is needed when stretching sible delimiters.
\[
\left. F(x) \right|_{a}^{b}\text{ note how this stretches the F to match the vertical bar}
\]

\section{Operators}
LaTeX defines e.g. trigonometric functions and mathematical limits as operators. The former does not involve limits while the latter can.
\begin{center}
	\begin{tabular}{llc}
	\toprule
	Typeset\\
	\midrule
	$\arccos$\\
	$\arcsin$\\
	$\arctan$\\
	$\arg$\\
	$\cos$\\
	$\cosh$\\
	$\cot$\\
	$\coth$\\
	$\csc$\\
	$\deg$\\
	$\dim$\\
	$\exp$\\
	$\hom$\\
	$\ker$\\
	$\lg$\\
	$\ln$\\
	$\log$\\
	$\sec$\\
	$\sin$\\
	$\sinh$\\
	$\tan$\\
	$\tanh$\\
	\bottomrule
	\end{tabular}
	\captionof{table}{Operators without limits (see source for type or command)}
	\label{Ta:opsnolimit}
\end{center}

\begin{center}
	\begin{tabular}{llc}
	\toprule
	Typeset\\
	\midrule
	$\det$\\
	$\gcd$\\
	$\inf$\\
	$\lim$\\
	$\liminf$\\
	$\injlim$\\
	$\varliminf$\\
	$\varinjlim$\\
	$\limsup$\\
	$\max$\\
	$\min$\\
	$\Pr$\\
	$\sup$\\
	$\projlim$\\
	$\varlimsup$\\
	$\varprojlim$\\
	\bottomrule
	\end{tabular}
	\captionof{table}{Operators with limits (see source for type or command)}
	\label{Ta:opsnolimit}
\end{center}

The limits to the operators are passed as parameters, for example:
\[
\projlim_{x \to 0}
\]

One can force limits in the subscript position with command \textbackslash nolimits:
\[
\projlim\nolimits_{x \to 0}
\]

\section{Multiline subscripts and superscripts}
It is also possible to apply multiline subscripts and superscripts, with \textbackslash substack. Each line is delimited with double backslashes:
\[
\sum_{\substack{i < n\\i \text{even}}}x_{i}^2
\]

To place the subscripts and superscripts flushed-left, use \textbackslash subarray instead of \textbackslash substack:
\[
\sum_{
\begin{subarray}{1}
i < n\\i \text{even}
\end{subarray}}
x_{i}^2
\]
\newpage
\section{Math accents}
Like text accents, these typically require a letter to accent but must be used within math environments.
\begin{center}
	\begin{tabular}{llc}
	\toprule
	Typeset\\
	\midrule
	$\acute{a}$\\
	$\bar{a}$\\
	$\breve{a}$\\
	$\check{a}$\\
	$\dot{a}$\\
	$\ddot{a}$\\
	$\dddot{a}$\\
	$\ddddot{a}$\\
	$\grave{a}$\\
	$\hat{a}$\\
	$\widehat{a}$\\
	$\mathring{a}$\\
	$\tilde{a}$\\
	$\widetilde{a}$\\
	$\vec{a}$\\
	\bottomrule
	\end{tabular}
	\captionof{table}{Math accents (see source for type or command)}
	\label{Ta:mathaccents}
\end{center}
\newpage
\section{Stretchable horizontal lines}
Three types in LaTeX: braces, bars and arrows.

Braces are typset with \textbackslash overbrace. A superscripts is optional. For example:
\[
\overbrace{a + b + \dots + n}^{m}
\]

Similarly, \textbackslash underbrace and optional subscripts are possible:
\[
\underbrace{a + b + \dots + n}_{m}
\]

Both braces can be nested.

Bars are typset with \textbackslash underline and \textbackslash overline as required, here shown as nested lines:
\[
\overline{
\overline{X} \cup \overline{\overline{X}}
} = \overline{\overline{X}}
\]

Arrows are either over or under and left, leftright or right, and have the same sort of features as braces and bars. Some examples:
\begin{gather*}
\overleftarrow{a}\\
\overrightarrow{ab}\\
\overleftrightarrow{abc}
\end{gather*}

The environment denoted by \texttt{gather} is used to gather single line formulae, in this case without numbering (as denoted by the asterisk).
\end{document}